\documentclass[11pt]{article}

\usepackage[T1]{fontenc}
\usepackage[utf8]{inputenc}
\usepackage{helvet}
\renewcommand{\familydefault}{\sfdefault}
\usepackage[a4paper, top=2.5cm, bottom=2.5cm, left=2.5cm, right=2.5cm]{geometry}
\usepackage{setspace}
\usepackage{parskip}
\usepackage{microtype}
\usepackage{amsmath}
\usepackage{booktabs}
\usepackage{hyperref}
\usepackage{verbatim}

\singlespacing

% ------------------------------------------------------------------
% AI tool declaration
% This document was drafted with assistance from Claude Sonnet 4.6
% (Anthropic, 2025) acting as a writing aid. All technical content,
% decisions, and code were produced by the authors.
% ------------------------------------------------------------------

\title{\textbf{Source3: A Rule-Based Pipeline for Equation Annotation\\
in Scientific Papers}}
\author{}
\date{}

\begin{document}

\maketitle

% ==================================================================
\section{Overview}
% ==================================================================

\subsection{Problem Analysis}

The goal of Source3 is to produce structured annotations for a curated
set of mathematical equations drawn from 46~arXiv papers in the domain
of quantum computing and quantum information.
For each reviewed equation the system must extract three types of
information:
\begin{enumerate}
  \item a \emph{meaning} --- a natural-language sentence that describes
        what the equation expresses;
  \item \emph{symbol definitions} --- a mapping from every mathematical
        symbol in the equation to its natural-language definition; and
  \item \emph{inter-equation relations} --- a classification of how the
        equation relates to every other reviewed equation in the same
        paper (e.g.\ derived from, equivalent, or unrelated).
\end{enumerate}
A central constraint is that all output strings must be verbatim spans
from the source paper; no text may be generated or paraphrased.
This requirement makes the outputs fully traceable and avoids
hallucination, but it forces every extraction step to locate the
correct span within the document rather than synthesising an answer.

Three technical challenges characterise the problem.
First, \emph{equation alignment}: the human-curated equation list
identifies each equation by its LaTeX source and, where available, by a
visible label or a DOM anchor extracted during a prior annotation pass.
The HTML rendered by arXiv uses a different identifier scheme and may
contain minor formatting differences, so a multi-strategy matching
procedure is required to link each reviewed equation to its unique DOM
node unambiguously.
Second, \emph{meaning localisation}: the sentence that best describes an
equation is not always the one immediately preceding it.
It may occur in a cross-reference later in the paper, in an equation
caption, or may be signalled by a cue phrase such as ``is defined as''
or ``Eq.~(N) gives''.
Naïve proximity heuristics are insufficient; ranking is needed.
Third, \emph{symbol disambiguation}: the same symbol (e.g.\ $\psi$) may
appear in many equations with distinct roles across a paper.
The correct definition must be retrieved for each specific occurrence
without confusion from other uses.

\subsection{Approach Overview}

Source3 is a fully deterministic, non-generative pipeline consisting of
ten sequential stages.
It uses a semantic embedding model (MathBERT~\cite{peng2021mathbert},
specifically the \texttt{witiko/mathberta} checkpoint) for dense
retrieval, BM25 for lexical matching, and hand-crafted linguistic
patterns for high-precision span extraction.
No large language model is invoked for generation.

The pipeline is divided into four logical phases.

\paragraph{Ingestion phase (Stages 1--4).}
Raw HTML is downloaded from arXiv and cached locally (Stage~1).
Each HTML document is then parsed into a hierarchical model comprising
sections, paragraphs, sentences with stable identifiers, and display
equation blocks with their associated MathML and original \LaTeX{} source
(Stage~2).
The document is subsequently segmented into typed retrieval chunks ---
individual sentences, paragraphs, raw equation neighbourhoods (up to five
sentences on each side of an equation), and cross-reference spans
(Stage~3).
Finally, every display equation is extracted from the parsed DOM: each
record captures its LaTeX source, anchor identifier, section, document
order, and the nearest prose sentences on each side (Stage~4).

\paragraph{Embedding phase (Stage 5).}
A single MathBERT instance is loaded once per paper.
Two families of vectors are computed and stored in a per-paper \texttt{.npz}
file.
Family~A is equation-centric: for each reviewed equation, up to twelve
vectors are produced encoding the equation alone, a summary context
(equation with five surrounding sentences), and individual
sentence-conditioned views (equation paired with each neighbouring
sentence).
Family~B is sentence-centric: every sentence and neighbourhood chunk in
the document is embedded to support symbol definition retrieval later.
All vectors are mean-pooled over the final hidden state, L2-normalised,
and stored as \texttt{float32[768]}.

\paragraph{Extraction phase (Stages 6--9).}
Equation meanings are retrieved by first collecting candidate sentences
from three sources in priority order --- cross-references, cue-pattern
matches, and proximity fallback --- then reranking survivors with a
weighted score combining cosine similarity to the equation's summary
vector, BM25, source-type bonus, and proximity (Stage~6).
A postprocessing pass refines or discards low-quality candidates
(Stage~6b).
Symbols are extracted from the MathML tree of each equation, with
composite identifiers (subscripted, superscripted, decorated) resolved
before flat leaf collection, and a LaTeX-tokeniser fallback applied when
MathML is absent (Stage~7).
Symbol definitions are located in two steps: an embedding pre-filter
computes cosine similarity between a query formed from the symbol's
aliases and all sentence-chunk vectors, retaining the top 30 candidates;
ordered linguistic patterns (``where \textit{sym} is'', ``\textit{sym}
denotes'', ``let \textit{sym} be'', etc.) are then applied to those
candidates only (Stage~8).
Inter-equation relations are scored using explicit citation cues,
derivation and equivalence cue phrases, Jaccard overlap of canonical
symbol sets, section proximity, and cosine similarity between summary
vectors; relations are graded as \textit{strong}, \textit{potential}, or
\textit{none} (Stage~9).

\paragraph{Export phase (Stage~10).}
The four intermediate artefacts (equations, meanings, symbol definitions,
relations) are joined on equation identifier, validated against a
checklist of schema invariants, and written atomically to
\texttt{data/final\_data.json} via a temporary file and rename.

\subsection{Pipeline Diagram}

The diagram below shows the data flow through the ten stages.

\begin{verbatim}
  [arXiv]
     |
     v
 Stage 1: Download HTML
     |  data/html/<paper_id>.html
     v
 Stage 2: Parse Documents
     |  documents/<paper_id>.json
     |  (sections, paragraphs, sentences, MathML)
     v
 Stage 3: Build Chunks
     |  chunks/<paper_id>.json
     |  (sentence | paragraph | neighborhood | cross_reference)
     v
 Stage 4: Extract Equations
     |  equations/<paper_id>.json
     |  (LaTeX, anchor_id, section, before/after sentences)
     v
 Stage 5: Build MathBERT Embeddings
     |  embeddings/<paper_id>.{npz,json}
     |  Family A: equation / summary / before_sentence / after_sentence
     |  Family B: sentence_chunk (whole document)
     |
     +------------------+------------------+------------------+
     |                  |                  |                  |
     v                  v                  v                  v
 Stage 6:           Stage 7:           Stage 8:           Stage 9:
 Equation           Symbol             Symbol             Relations
 Meanings           Extraction         Definitions        Detection
 (cosine+BM25       (MathML tree       (embed pre-        (cue phrases,
  + cue patterns)    traversal)         filter + patterns)  symbol overlap,
     |                  |                  |                semantic sim.)
     v                  v                  v                  v
 meanings/          symbols/           symbol_meanings/   relations/
 <paper_id>.json    <paper_id>.json    <paper_id>.json    <paper_id>.json
     |                  |                  |                  |
     +------------------+------------------+------------------+
                        |
                        v
                   Stage 10: Export
                        |
                        v
               data/final_data.json
\end{verbatim}

\begin{thebibliography}{9}
\bibitem{peng2021mathbert}
S.~Peng, K.~Yuan, L.~Gao, and Z.~Tang,
``MathBERT: A Pre-trained Model for Mathematical Formula Understanding,''
\textit{arXiv preprint arXiv:2105.00377}, 2021.
\end{thebibliography}

\end{document}
